\subsection{Gene dependence in factor--phenotype enrichment}
\label{sec:gene_dependence_robustness}

\paragraph{Statistical question and provenance of significance.}
For a fixed factor and target phenotype, gene-level enrichment asks whether factor membership predicts the target response beyond a fitted baseline and any explicitly included covariates. Its null permits a nonzero overall hit rate. It does not require the target to have no genetic associations. An intercept accounts for a fitted mean; it does not remove covariance between gene responses. Moreover, reference inclusion and factor/trait selection can change the comparison population. In particular, a Direct-score analysis restricted to a Combined-selected reference is not an uncensored Direct-only analysis.

Three computational paths must be distinguished. The public fixed-basis NNLS linkage is descriptive and has no association P value. In native EAGGL binary enrichment, the implemented model is a linear probability regression, with independent-observation HC3 uncertainty by default. HC3 allows heterogeneous residual variance but contains no off-diagonal gene covariance. The current native binary fitter does not implement chromosome-clustered uncertainty. Separately, exported factor GMTs can be tested using PIGEAN multi-Y logistic regression. There, the output \texttt{P}, \texttt{Z}, and \texttt{SE} accompany the marginal \texttt{beta\_tilde} statistics. The downstream Bayesian \texttt{beta} and \texttt{beta\_uncorrected} columns are different estimates; their presence does not turn that P value into a simultaneous all-factor test. The dashboard carries these fields alongside each other. A regression test verifies that changing the Bayesian beta does not change the imported marginal P value.

The audited T2D factor-GMT run selected Combined responses, background prior 0.05, logistic rather than linear regression, and no gene-correlation or gene-location input. Its independent-gene uncertainty did not invoke the annotation-derived continuous-response correction. Optional gene covariate corrections and Bayesian factor adjustment must likewise not be described as an annotation-dependence correction. The subsequent explicit residual-correlation interface corrects marginal statistics; it is not a newly derived covariance-aware joint logistic likelihood.

\paragraph{What the covariance correction establishes.}
For a marginal logistic fit, let $v_i=\widehat p_i(1-\widehat p_i)$ and
\[
 c_i=\sqrt{v_i}\left(x_i-\frac{\sum_jv_jx_j}{\sum_jv_j}\right).
\]
For a specified residual correlation matrix $R$, intercept-adjusted uncertainty is
\begin{equation}
 \mathrm{SE}_R=\mathrm{SE}_I
 \sqrt{\frac{c^\top R c}{c^\top c}}.
 \label{eq:gene_dependence_sandwich}
\end{equation}
The existing independent synthetic-pseudocount convention is included in both quadratic forms. Identity and independently assembled dense-sandwich tests verify this calculation. For linear regression, $c_i=x_i-\bar x$; the simultaneous linear covariance additionally includes the fitted residual-variance scale. These checks establish the algebra for a specified matrix, not that a matrix describes the biological pipeline.

A shared-annotation working matrix can represent similarities in annotation contributions, including similarities between genes on different chromosomes. It does not automatically propagate uncertainty in the learned annotation coefficients or gene-reference selection. Treating Direct support as independent diagonal variance is also a modeling assumption: physical linkage can correlate Direct evidence. Genomic clustering addresses a different dependence structure and, by itself, cannot address shared annotations across chromosomes. Finally, thresholding continuous Combined support changes the covariance; positive semidefiniteness and unit diagonal alone do not establish a valid binary-response correlation model.

For two Bernoulli responses with means $p_i,p_j$, a necessary pairwise bound is
\begin{equation}
 \frac{\max(0,p_i+p_j-1)-p_ip_j}{\sqrt{p_i(1-p_i)p_j(1-p_j)}}
 \le R_{ij}\le
 \frac{\min(p_i,p_j)-p_ip_j}{\sqrt{p_i(1-p_i)p_j(1-p_j)}}.
 \label{eq:binary_correlation_bounds}
\end{equation}
A working matrix can fail this bound even while being positive semidefinite. Passing pairwise bounds is necessary but is not sufficient to establish a compatible multivariate binary distribution or empirical calibration.

\paragraph{Deterministic HDL sensitivity audit.}
We retained the saved 18,321-gene universe, the Factor 4 and Factor 7 GMT memberships, and the 863 HDL Combined-hit labels. Factor 4 contained 323 genes and 74 hits; Factor 7 contained 173 genes and 87 hits. The fitted response labels were held fixed throughout deletions. Closed-form two-group logistic fits agreed with the production fitter at tighter numerical convergence tolerance; minor differences from earlier rounded/default-tolerance statistics do not affect interpretation. The table below uses the corrected marginal fitter, rather than the historical dashboard export that preceded the trait-batching correction.

The saved reference provided HDL gene scores but not HDL's fitted annotation-coefficient table. We therefore did not claim to reconstruct its trait-specific annotation covariance. Instead, a target-independent structural-overlap stress used the four supplied reference annotation sources, with 81,452 unique gene-membership sets after removing 6,544 duplicates. This basis covers 99.28\% of the gene universe and is broader than the selected T2D discovery annotations. Let $B$ be its binary gene-by-annotation matrix, and let $L$ normalize each nonzero gene row to unit Euclidean norm. Define
\begin{equation}
 K=LL^\top+\operatorname{diag}\{\mathbf1[\|B_{i\cdot}\|=0]\},
 \qquad R_\lambda=(1-\lambda)I+\lambda K.
 \label{eq:annotation_overlap_stress}
\end{equation}
This matrix is computed through sparse products without allocating a dense gene-by-gene matrix. We examined the fixed grid $\lambda\in\{0,0.1,0.25,0.5,0.75\}$. These are declared dependence assumptions, not estimates of residual correlation. The tests do not use target association strength to select annotations or tune $\lambda$.

\begin{center}
\begin{tabular}{lrrr}
Factor & $\lambda$ & SE multiplier & Working-model P value\\
\hline
4 & 0 & 1.000 & $4.18\times10^{-42}$\\
4 & 0.10 & 1.314 & $4.15\times10^{-25}$\\
4 & 0.25 & 1.677 & $5.25\times10^{-16}$\\
7 & 0 & 1.000 & $2.08\times10^{-88}$\\
7 & 0.10 & 1.260 & $2.37\times10^{-56}$\\
7 & 0.25 & 1.572 & $7.45\times10^{-37}$\\
\hline
\end{tabular}
\end{center}
The displayed settings pass the necessary pairwise binary bounds under each fitted model. Settings 0.5 and 0.75 fail those bounds for both factors; their numerical results are retained in the diagnostic tables but are not presented as valid binary-model inference. At $\lambda=0.1$, omitting one annotation source at a time gives Factor 7 SE multipliers 1.081--1.286 and working-model P values from $7.22\times10^{-76}$ to $3.23\times10^{-54}$. Exact duplicate memberships supported by another retained source remain represented.

We separately removed genes in each of the ten annotation groups with largest overlap with each fixed factor, ranked without HDL evidence, and every supplied 1-Mb gene-start window overlapping that factor. The latter are coordinate-based influence groups, not validated LD blocks; genes with ambiguous mappings enter each applicable deletion. Across 171 Factor 7 windows, all remaining log-odds coefficients exceeded 3.084, compared with 3.110 initially. Across its ten annotation-group deletions, all exceeded 2.832. The largest coefficient reduction followed deletion of the oxygen-containing-compound response annotation, leaving 51.4\% of factor members. Factor 4 coefficients remained above 1.839 across 291 windows and above 1.392 across ten annotation deletions. No fit became unidentified or separated. These are fixed-label influence checks, not refits of annotation learning; they do not prove that removing an annotation from PIGEAN would leave the same gene-hit list.

\paragraph{Interpretive limit.}
The positive enrichment is not explained by a single tested genomic window or major overlapping annotation group, and it persists under the specified feasible-pairwise-bound correlation stresses. The precise magnitude of its P value is sensitive to dependence assumptions. These observations support robustness of the observed enrichment, not calibration of the extreme independent-gene P value, exhaustive coverage of dependence structures, or confirmation in an independently selected factor/trait panel. No new simulations were undertaken for this deterministic audit. Additional calibration would be needed if a claim depended on precise tail probabilities that these checks cannot establish.
